\documentclass[11pt]{letter}
\usepackage[T1]{fontenc}
\usepackage{lmodern}
\usepackage[margin=1in]{geometry}
\usepackage{microtype}
\usepackage{xcolor}
\usepackage{enumitem}
\setlist{itemsep=2pt,topsep=4pt}
\newcommand{\todo}[1]{\textcolor{red!70!black}{[TODO: #1]}}
\newcommand{\MsTitle}{Neighborhood faithfulness in maps of trial-averaged neural responses, estimated from held-out repetitions}
% The numbers the letter quotes are written by make_tables.py from the committed results,
% between its BEGIN and END lines; python3 make_tables.py --check verifies them.
% BEGIN make_tables.py numbers
\newcommand{\EegLeight}{64}
\newcommand{\EegFCeightSmaller}{81}
\newcommand{\BmEnrichLo}{2.0}
\newcommand{\BmEnrichHi}{4.2}
\newcommand{\BmFiveLo}{68}
\newcommand{\BmFiveHi}{73}
% END make_tables.py numbers
\signature{Sooyoung Jang\\ Hanbat National University\\ \texttt{syjang@hanbat.ac.kr}}
\address{Sooyoung Jang\\ \todo{department}\\ Hanbat National University\\ 125 Dongseo-daero, Yuseong-gu\\ Daejeon 34158, Republic of Korea}

\begin{document}
\begin{letter}{Professor Giuseppe Di Giovanni\\ Editor-in-Chief\\ Journal of Neuroscience Methods}
\opening{Dear Professor Di Giovanni,}

Please consider the enclosed research article, ``\MsTitle'', for publication in the
\emph{Journal of Neuroscience Methods}.

Two-dimensional maps of trial-averaged responses, made with t-SNE or multidimensional scaling, are
read neighbor by neighbor throughout systems neuroscience, and analysts check them against the
averaged data or for stability across runs, resamples and halves of the trials. The manuscript
asks how often a neighborhood that passes such a check is faithful to the noise-free responses,
and answers with a method that estimates this from held-out repetitions and says when its estimate
can be trusted.

The study reanalyzes two public datasets, THINGS-EEG2 and the Allen Brain Observatory two-photon
survey, and I am aware of the journal's policy for such work. I believe it meets that policy on
all three of its grounds:
\begin{itemize}[leftmargin=*]
  \item \textbf{Methodological innovation.} An estimator of the ``false comfort'' of any
  neighborhood check, from held-out reference blocks, with an extrapolation for the reference's
  own noise and subject-level bootstrap bounds; and a simulation study built on the real data that
  defines observable statuses, saying when the estimate can be trusted, and tests whether its
  bound is calibrated across simulated studies.
  \item \textbf{Rigorous benchmarking.} A comparison of six existing checks against the truth and
  against a random selection of as many neighborhoods: per-point trustworthiness, scDEED,
  cross-validated distances, and recurrence across t-SNE seeds, bootstrap resamples and split
  halves. In EEG the checks selected faithful neighborhoods \BmEnrichLo--\BmEnrichHi{} times more
  often than chance, yet even their best-scoring 5\% had an estimated false comfort of
  \BmFiveLo--\BmFiveHi\%.
  \item \textbf{Insight for experimental methodology.} At 8 trials, an estimated
  \EegFCeightSmaller\% of the neighborhoods that pass the usual check are not faithful in human
  EEG, with a lower bound of \EegLeight\% that is calibrated in simulated studies. In mouse V1, a
  held-out reference of 24 trials confirmed fewer than 0.1\% of them. The statuses also show that
  at 16 or 32 trials in EEG, and in mouse V1, held-out references of the depth these designs allow
  cannot resolve false comfort, which bears directly on how many repetitions a study needs before
  its maps can be read locally.
\end{itemize}
Every claim, threshold and robustness variation was fixed in a protocol committed before the
analysis it governs was run, and the mouse data were selected by rule and analyzed only after
their protocol was fixed. The simulation study and the rules for reporting each result were
specified after the first results, in a further protocol committed before they were run; the
manuscript reports each result according to the status those rules give it. In line with the
journal's research data policy, the code, protocols and derived results are deposited at
\todo{repository and DOI} and cited in the article; the raw data remain openly available from
their providers.

The manuscript has not been published and is not under consideration elsewhere. The author
declares no competing interests. \todo{confirm} As described in the manuscript's Methods and its
declaration of generative AI use, an AI assistant (Claude, Anthropic) was used under the author's
direction to design, implement and run the analyses and to draft the text, figures and tables; the
author reviewed and edited the content and takes full responsibility for it.

\closing{Sincerely,}
\end{letter}
\end{document}
